\section{Discussion}

\textit{Written 2026-07-05. Interprets \texttt{paper\_results.md} R1--R6; numbers there are canonical.}

\subsection{1. What the results mean together}

The five analyses tell one coherent story that refines --- and in one place inverts --- the field's prevailing intuition.

\textbf{The unifying result is the five-setting map, not any single pairwise comparison (R6e).} Ordering NTU$\rightarrow$\{CZU skeleton, UTD skeleton, Xsens quats, CZU IMU orientation-only, CZU IMU dual-raw\} by gap width and target-representation strength, both the \textit{best} prior and the \textit{actively harmful} prior change identity as the setting moves: supervision wins at the two small, same-modality gaps; the hybrid wins and pure reconstruction is the harmful one at the middle, cross-device gap; nothing beats scratch and pure supervision is the harmful one at the two large, cross-modal gaps. This "culprit swap" --- mae is the negative-transfer case only at the middle setting (R2), supLP120 only at the two large-gap settings (R6b/R6c) --- is not an inconsistency to explain away; it is the finding. No pretraining objective is unconditionally safe across the skeleton$\rightarrow$wearable gap, and \textit{which} objective is unsafe is itself gap-dependent. The analyses below (R1--R6d) are what establish and stress-test this map; we walk through them in turn.

\textbf{Objective choice is real but gap-gated (R1).} The 45 pp within-domain spread collapsing to \textasciitilde{}4 pp at k=1 across the modality gap is not evidence that objectives are interchangeable; it is evidence that a large representation gap acts as a low-pass filter on whatever structure the source objective built. This reconciles the two camps in the literature: the "objectives matter" results come from small-gap or within-modality settings, and the "just use a strong SSL objective" results come from settings where the gap has already been narrowed by scale or modality proximity.

\textbf{The negative-transfer case inverts "reconstruction is safe" --- at the middle-sized gap (R2).} The field has drifted from supervised to reconstruction/hybrid objectives on the premise that generative, structure-preserving objectives are the robust default under shift. Our controlled cross-modal result complicates that premise: \textbf{pure masked reconstruction (mae) is the significant negative-transfer case at every shot count} (McNemar p<.001), while the \textit{hybrid} is robustly best and \textit{pure supervision} is competitive and best-calibrated. The resolution is that "richer objective" was always a proxy for "objective that preserves discriminative structure while remaining transferable." MAE alone satisfies neither in the cross-modal regime: it produces a barely-discriminative representation within domain (22.6\%), and starting the fine-tune from that biased init lands in a worse basin than random. The hybrid keeps the discriminative signal \textit{and} the reconstruction regularizer, which is why it, not pure reconstruction, is the safe choice here.

\textbf{Alignment is necessary, not sufficient (R3).} CKA cleanly orders the objectives (supervised aligns best, reconstruction barely above random) and does so monotonically with depth --- a satisfying, scale-invariant mechanism the MMD-only view could not provide. But the best-aligned encoder is not the best-transferring, and the reconstruction encoder aligns better than scratch yet transfers worse. Representation similarity sets an upper bound on what \textit{can} transfer; the objective's discriminative content decides what \textit{does}. Stating this two-part conclusion --- gap necessary, discriminative structure sufficient --- is more honest and more useful than a single-metric claim.

\textbf{Encoder alignment and raw distributional gap turn out to be different quantities (R3).} When we extended CKA across all five source$\rightarrow$target settings to test whether it tracks the paper's claimed small/middle/large gap ordering, it did not --- czu\_skeleton (claimed small gap) ranked as the \textit{most} CKA-dissimilar of the four targets on the best-aligning encoder, below even the claimed-large setting. Rather than let that stand as an unresolved crack in the gap ordering, we measured the gap a second way, directly on a hand-crafted, encoder-free feature (no pretraining, no learned representation at all): squared-MMD and Frechet distance between NTU and each target. Both confirm the small<small<middle<large ordering exactly and monotonically. The two results are not in tension once read correctly: CKA measures how well a \textit{specific trained encoder's} representation lines up across domains, which is objective-dependent by construction (that dependence is R3's main finding); raw MMD/Frechet measure how different the \textit{datasets themselves} are before any model touches them. A metric computed inside a trained encoder inherits that encoder's objective-specific biases as a confound on any claim about the datasets, which is exactly why it failed to track gap width here. The five-setting map's gap ordering now rests on this encoder-free measurement, not narrative alone.

\textbf{The prior's value is calibration and data-efficiency, not headline accuracy (R4).} This is the most deployment-relevant reframe. Where endpoint accuracy is a wash, the prior still calibrates better (supLP120 lowest ECE at every k), integrates few shots faster (AUC-30 and convergence significant), and --- most concretely --- pays off most exactly when target data is scarcest: its benefit peaks at 0--1 enrolled subjects (+27 pp) and washes out by four. For a practitioner the actionable question is not "which objective wins at full data" (none, materially) but "how many users must I enroll before the cross-domain prior stops earning its keep" --- and the answer, in this setting, is about four.

\textbf{Small differences compound into task consequences (R5).} The controller shows that a 4 pp recognition gap becomes a 21 pp task-success gap under compounding, and that calibration quality --- not just accuracy --- sets the safety/throughput frontier. This is the human--machine-systems payload: it converts a within-noise accuracy table into a decision-relevant reliability story, and it does so on the \textit{actual} posterior stream rather than a synthetic error model.

\textbf{The cross-modal signature is specific, and prior value is gap-contingent (R6/R6b/R6c/R6d, all pooled over 3 seeds).} That mae hurts on NTU$\rightarrow$Xsens but \textit{helps} on CZU skeleton$\rightarrow$skeleton is a control, not a footnote: it shows the negative transfer is a property of the \textit{modality gap}, not of the mae objective in the abstract, exactly as the mechanism (R3) predicts. Four independently pooled legs across three public datasets sharpen this into a monotone arc, and multi-seed pooling changed \textit{which} claim survives at two of those legs. Same-modality, the pure-supervised prior supLP120 is best-in-class and now clearly significant rather than marginal: \textbf{+7.2 pp zero-shot on CZU (p=.0004, up from single-seed p=.07)}, independently replicated on a second dataset, \textbf{+7.4 pp on UTD-MHAD (p<.0001)} --- two datasets, not one, agree that supervision wins when the gap is small. Across the modality gap with an impoverished orientation-only target (R6b), the ordering inverts, but pooling corrected the headline: the single-seed "supMAE beats scratch by 3 pp" did \textbf{not} survive (25/42 folds, p=.28 --- a seed-42 artifact); what is robust is the \textit{prior-vs-prior} ordering, \textbf{supMAE > supLP120 in 36/44 folds (p<.0001)}, with supLP120 below scratch at every k (p$\approx$.04--.06). The honest R6b claim is therefore supervised-specific negative transfer neutralized by reconstruction back to $\approx$scratch, not reconstruction beating scratch outright. And once the target is given its full raw signal (R6c dual-branch, where from-scratch already matches CRC), even that neutralization vanishes: supMAE ties scratch (17/39, p=.52, no detectable benefit) and supLP120 is again worst, now more decisively (6/41, \textbf{p<.0001}, paired-t significant at every k). The \textit{same} prior thus moves hero (two datasets, small gap) $\rightarrow$ villain (large gap, weak target) $\rightarrow$ still villain (large gap, strong target), and the R6b reconstruction rescue is revealed as contingent on target-representation poverty rather than a general property of cross-modal transfer. A transferred prior earns its keep only in proportion to how impoverished the target representation is --- which is the negative-transfer thesis stated as a design rule, now demonstrated on three independent public datasets rather than asserted on one.

\subsection{2. Relation to prior work}

Our within-IMU reference points (BenchHAR, PRIMUS) find hybrid recon+contrastive best under OOD shift; we corroborate "hybrid is robust" in the harder \textit{cross-modal} skeleton$\rightarrow$IMU setting, but add the sharper finding that \textit{pure} reconstruction can be actively harmful there --- a distinction those within-modality studies could not surface because their gap is small. Against the supervised-transfer line (Moya Rueda \& Fink, Awasthi), which reports supervised gains under moderate gaps, we neither confirm a supervised win nor the assumed supervised \textit{collapse}: supervision is gap-gated to near-neutral on accuracy but retains the best calibration and the strongest low-enrollment benefit. Against the foundation/contrastive line (IMU2CLIP, UniMTS, AURA-MFM), we do not compete on scale; our contrastive arm (SupCon) is a scoped negative exhibit under one preprocessing, and we are explicit that controlled small-N objective isolation is a different question than foundation-scale alignment.

\subsection{3. Limitations (owned)}

\begin{itemize}
  \item \textbf{Five subjects, one primary dataset pair.} We address power with 3-seed pooling (n=15) and clip-level McNemar (\textasciitilde{}400--500 paired clips/subject) rather than the n=5 Wilcoxon floor (p$\ge$0.0625), and we frame the work as a controlled characterization, not a benchmark. External validity is discharged across four pooled, multi-seed legs on three public datasets --- same-modality CZU (R6, n=15) and UTD (R6d, n=24), cross-modal orientation-only CZU (R6b, \textasciitilde{}44 folds), and cross-modal dual-branch CZU (R6c, 45 folds) --- so the cross-modal effects rest on well-powered per-fold sign tests (e.g. supMAE > supLP120 36/44 p<.0001 in R6b; supLP120 worst 6/41 p<.0001 in R6c), not the single-seed n=5 floor. Two pieces remain single-seed: the R6c raw-only/quat-only diagnostic rows (the pooled \textit{dual} mode is what carries the claim) and the target-richness dose-response dial. The R6c supMAE$\approx$scratch and R6b supMAE$\approx$scratch results are nulls at this power (no \textit{detectable} benefit, not proven equivalence) --- pooling in fact reversed the single-seed R6b headline (supMAE beating scratch was a seed-42 artifact, 25/42 p=.28), which is itself evidence the pooling was necessary rather than cosmetic.
  \item \textbf{Contrastive is under-sampled.} SupCon was run only under the earlier single-seed local preprocessing, where it underperformed scratch. We report it as a scoped exhibit, not a multi-seed v2 arm; a full contrastive sweep is future work and we do not draw strong contrastive conclusions.
  \item \textbf{The controller's magnitudes are cost-model-dependent --- but its ordering is not.} Absolute task-success numbers depend on the System Input assignment, the error-cost model, and the operating threshold. Rather than freeze these, we show (Results R5) that the two claims that matter --- \textit{mae compounds worst} and \textit{calibration governs the safety/throughput trade} --- survive across 120 random System Input assignments, two outcome models with a 6-point critical-cost sweep, and a tuning-free iso-safety operating point. We therefore present the controller as a demonstration that objective/calibration differences propagate to task reliability with the right sign and ordering, not as a calibrated estimate of any specific deployment's success rate. One nuance is itself a finding: the best-calibrated init (supLP120) has a confident false-critical-activation mode that penalizes it on the \textit{ungated} metric, which is exactly why the deployment-correct framing is iso-safety plus the design guard of assigning the Safety-Critical States to the most-separable gestures. Inter-command timing is synthesized.
  \item \textbf{Subject-scaling is single-seed.} The A2 k=3 trend is robust; the k=1, N=1 spike sits in the seed-sensitive k=1 band and should be read as suggestive.
  \item \textbf{The swing preprocessing is a methodological finding, not an experimental lever.} We report it as evidence that surface preprocessing and feature-distance metrics are necessary-not-sufficient, having falsified its original gap-closing hypothesis; it is not part of any headline claim.
\end{itemize}

\subsection{4. Practical guidance}

For a practitioner building a wearable gesture recognizer with a skeleton prior: (i) prefer a \textbf{hybrid reconstruction+supervision} objective --- pure reconstruction can hurt across a modality gap; (ii) expect the prior's accuracy benefit to be \textbf{small at full data but large at 0--1 enrolled users}, and to manifest as \textbf{better calibration and faster calibration convergence} even when endpoint accuracy is a wash; (iii) if the downstream task is a \textbf{safety-sensitive controller}, the calibration quality of the objective matters as much as its accuracy, because confidence-threshold operating points --- not raw accuracy --- set the safety/throughput trade; and (iv) \textbf{weigh the prior against the target's own signal} --- R6c shows that once the target model can exploit its full raw stream (matching a hand-crafted baseline from scratch), the skeleton prior adds no accuracy and a pure-supervised prior can subtract it, so the prior is worth most exactly where the target representation is impoverished.
